\section{Apéndice III: tarjeta de autoevaluación para laboratorios de modelos}

Para convertir esta entrevista en criterios aplicables, puede elaborarse una tarjeta de evaluación simplificada para cualquier laboratorio de modelos o equipo emprendedor de Agent. No es una herramienta de clasificación, sino una ayuda para que el lector juzgue si un equipo comprende de verdad el paradigma Agent.

\noindent\begin{longtable}{p{0.24\textwidth}p{0.33\textwidth}p{0.34\textwidth}}
\toprule
Dimensión & Desempeño de puntaje alto & Señales de puntaje bajo \\
\midrule
Juicio sobre el paradigma & Explica con claridad las diferencias de entrenamiento y de producto entre las tres etapas: Chat, Reasoning y Agent & Concibe Agent solo como llamadas a herramientas o como un prompt complejo. \\
Sistema de posentrenamiento & Cuenta con un ciclo cerrado de entornos de tareas, rollout, verifier, eval y monitoreo de costos & Solo dispone de datos fuera de línea o de ejemplos seleccionados manualmente. \\
Orquestación de modelos & Enruta modelos grandes y pequeños, modelos multimodales y modelos especializados según la etapa de la tarea & Todas las tareas llaman al mismo modelo, el más grande. \\
Asignación de cómputo & research, pre-train y post-train tienen proporciones claras y un ritmo de iteración definido & Los recursos quedan atados a la trayectoria histórica y el posentrenamiento no tiene suficientes GPU. \\
Mecanismos organizacionales & Preentrenamiento, posentrenamiento, producto, sistemas y evaluación participan juntos en las decisiones & Un solo equipo decide a puerta cerrada; la retroalimentación de usuarios reales y de ingeniería no llega. \\
Conciencia de costos & Optimiza en conjunto la tasa de finalización de extremo a extremo, la latencia y el precio & Solo le interesa exhibir capacidades, no la economía por tarea. \\
Valores y límites & Define con claridad qué tareas deben automatizarse y en cuáles las personas deben conservar el juicio & Solo persigue la sustitución y el crecimiento, sin discutir seguridad, responsabilidad ni consecuencias sociales. \\
\bottomrule
\end{longtable}

\begin{knowledgebox}{Cómo usar esta tarjeta de evaluación}
Al leer las siguientes entregas de la serie Zhang Xiaojun AI, pueden situarse las opiniones de cada invitado en estas siete dimensiones. Por ejemplo, la entrevista con Su Yu (苏煜) se inclina más hacia «juicio sobre el paradigma» y «aprendizaje a largo plazo», mientras que la entrevista con Luo Fuli (罗福莉) se inclina más hacia «sistema de posentrenamiento», «orquestación de modelos», «asignación de cómputo» y «mecanismos organizacionales». Así, las notas de varias entregas irán formando un mapa comparable del sistema operativo de un AI Lab.
\end{knowledgebox}

\subsection{Resumen de la sección}

La tarjeta de evaluación convierte los juicios de la entrevista en un marco reutilizable. Recuerda al lector que, en la era de los Agent, no triunfa una sola tecnología aislada: varias dimensiones del sistema deben cumplir el estándar al mismo tiempo.
